\section{Verifier 与 LLM-as-a-Judge：把评估嵌入推理}

\subsection{从终局评估到中间评估}
如果只在轨迹末尾判定成功，反馈太晚。讲者提出把 verifier 前移到中间步骤，让系统在推理时就能判断 ``当前路径是否值得继续''。

\begin{importantbox}{Verifier 的三类作用}
\begin{enumerate}
    \item 路径打分：为搜索裁枝提供依据；
    \item 错误检测：识别明显偏航状态；
    \item 数据过滤：筛选可用于训练的高质量轨迹。
\end{enumerate}
\end{importantbox}

\subsection{LLM-as-a-Judge 的现实价值}
让 LLM 评估轨迹通常比让它直接生成最优轨迹更稳定。因为 ``判断好坏'' 通常比 ``一步到位做对'' 更容易，这也是本讲反复强调的工程经验。

\begin{knowledgebox}{Judge 提示词的关键要素}
\begin{itemize}
    \item 任务目标与约束；
    \item 完整轨迹（或关键步骤摘要）；
    \item 明确输出模板：通过/失败 + 依据；
    \item 置信度估计，供搜索策略使用。
\end{itemize}
\end{knowledgebox}

\subsection{Judge 可靠性与校准}
Judge 不是绝对真值。若评估标准模糊或提示词漂移，Judge 可能给出高置信度错误结论，反而误导搜索。

\begin{warningbox}{Judge 的三种风险}
\begin{itemize}
    \item 奖励黑客：模型学会迎合 Judge 文风而非真实完成任务；
    \item 领域偏差：Judge 在陌生网页模板下失真；
    \item 置信度失配：高分不代表高成功率。
\end{itemize}
\end{warningbox}

\begin{figure}[H]
\centering
\includegraphics[width=0.82\textwidth]{frame-08-full-pipeline.jpg}
\caption{视频约 01:12:10：规划、执行、验证器构成完整闭环，评估不再只发生在任务末尾}
\end{figure}

\subsection{本章小结}
Verifier 让 Agent 从 ``先做再判'' 升级到 ``边做边判''。但 Judge 可靠性需要持续校准，否则会把系统引向错误优化方向。

